\subsection{有限支集分布}

13.6.1. 设 X 是一个 \(\mathbf{C}^r\) 类流形，并设 \(k \leqslant r\)。用 \(\mathcal{T}^{(k)}(\mathbf{X})\) 表示诸空间 \(\mathrm{T}_x^{(k)}(\mathbf{X})\)（\(x \in \mathbf{X}\)）的直和。\(\mathcal{T}^{(k)}(\mathbf{X})\) 的元素称为 \(\mathbf{X}\) 上阶 \(\leqslant k\) 的有限支集分布。若 \(f: \mathbf{X} \to \mathbf{F}\) 是 \(\mathbf{C}^r\) 类、取值于分离的多范数空间 \(\mathbf{F}\) 的函数，且 \(t = \sum_{x \in \mathbf{X}} t_x\)（其中 \(t_x \in \mathrm{T}_x^{(k)}(\mathbf{X})\) 对一切 \(x \in \mathbf{X}\) 成立）是一个有限支集分布，则我们令

\[
\langle f, t \rangle = \sum_ {x \in X} \langle f, t _ {x} \rangle .
\]

类似地，我们令 \(c(t) = \sum_{x \in X} c(t_x)\)，这给 \(\mathcal{T}^{(k)}(X)\) 装备了一个以 \(t \mapsto \langle 1, t \rangle\) 为余单位的余结合、余交换余代数结构。包含映射 \(\mathcal{T}^{(k')}(X) \to \mathcal{T}^{(k)}(X)\)（其中 \(k' \leqslant k\)）是余代数的态射。

若 \(\varphi: X \to Y\) 是一个 \(C^r\) 类流形态射，则诸映射 \(\mathrm{T}_x^{(k)}(\varphi): \mathrm{T}_x^{(k)}(X) \to \mathrm{T}_{\varphi(x)}^{(k)}(Y)\)（其中 \(x\) 取遍 \(X\)）定义了一个线性映射 \(\varphi_*\)，它把 \(\mathcal{T}^{(k)}(X)\) 映到 \(\mathcal{T}^{(k)}(Y)\)，且是余代数的态射。

若 \(X_{1}\) 与 \(X_{2}\) 是两个 \(C^{r}\) 类流形，则诸映射 \(\alpha_{k}^{-1}\)（参见 13.4.5）定义了一个线性映射

\[
\mathcal {T} ^ {(k)} \left(\mathrm{X} _ {1} \times \mathrm{X} _ {2}\right)\rightarrow \mathcal {T} ^ {(k)} \left(\mathrm{X} _ {1}\right) \otimes \mathcal {T} ^ {(k)} \left(\mathrm{X} _ {2}\right)
\]

它是单射；它是余代数的态射。若 \(k = \infty\) 或 \(\omega\)，它是余代数的同构。

13.6.2. 设 X 是一个 \(C^{r}\) 类流形，V 是一个有限维 K-向量空间。\(\mathcal{T}^{(r)}(\mathrm{X})\otimes\mathrm{V}\) 的元素称为 X 上取值于 V 的有限支集分布。当 K = R、V = C 时，这样的分布也称为 X 上的复有限支集分布。前面各条的定义和结论通过线性性立即推广到取值于 V 的分布。

